% Full answers PDF: MCQ answers + worked explanations per option, then official MS clips.
% Compiler: XeLaTeX.
\documentclass[12pt,a4paper]{article}
\usepackage[margin=18mm,headheight=14pt]{geometry}
\usepackage{fontspec}
\usepackage[version=4]{mhchem}
\usepackage{graphicx}
\usepackage{xcolor}
\usepackage{booktabs}
\usepackage{fancyhdr}
\usepackage{enumitem}
\usepackage{needspace}

\setmainfont{Helvetica Neue}

\pagestyle{fancy}
\fancyhf{}
\fancyhead[L]{\small 3.4 sigma/pi and bond energy --- answers}
\fancyhead[R]{\small CIE 9701}
\fancyfoot[C]{\thepage}
\renewcommand{\headrulewidth}{0.3pt}

\setlength{\parindent}{0pt}
\setlength{\parskip}{0.35em}

\newcommand{\msblock}[2]{%
  \par\vspace{0.8em}%
  \noindent\begin{minipage}{\textwidth}%
  {\normalsize\bfseries #1}\par\vspace{0.25em}%
  \noindent\includegraphics[width=\textwidth,keepaspectratio]{#2}\par
  \end{minipage}\par
}

\newcommand{\mcq}[4]{%
  \par\needspace{5\baselineskip}\vspace{0.7em}%
  {\normalsize\bfseries #1.\quad #2}%
  \hfill{\footnotesize\color{gray}#3}\par\vspace{0.15em}%
  {\small #4}\par
}

\newcommand{\why}[2]{\par\needspace{2.2\baselineskip}\vspace{0.1em}{\small\textbf{#1} #2}\par}

\begin{document}

\begin{center}
{\large\bfseries 3.4 Paper 2 $\sigma/\pi$ and bond energy homework --- answers}\\[0.2em]
{\small CIE 9701 AS Chemistry \quad·\quad 8 multiple-choice + 12 written marks}
\end{center}

{\small\bfseries Do not issue this file to students.}
\hrule
\vspace{0.4em}

\section*{Section A --- Multiple choice}

{\normalsize\bfseries Quick answers}

\vspace{0.3em}
{\small
\begin{tabular}{@{}ll@{\hspace{2.6em}}ll@{\hspace{2.6em}}ll@{}}
\toprule
Q & Ans. & Q & Ans. & Q & Ans.\\
\midrule
A1 & \textbf{B} & A4 & \textbf{C} & A7 & \textbf{C}\\
A2 & \textbf{A} & A5 & \textbf{D} & A8 & \textbf{B}\\
A3 & \textbf{C} & A6 & \textbf{B} & & \\
\bottomrule
\end{tabular}}

\vspace{0.6em}
{\normalsize\bfseries Worked explanations}

\mcq{A1}{B}{2025 ON Paper 12 Q5}{%
In \ce{N2} each nitrogen atom is \emph{sp} hybridised (one \ce{2s} mixes with one
\ce{2p}). A $\sigma$ bond is always formed by \emph{head-on} overlap, here
\emph{sp}--\emph{sp} overlap. Each N uses one \emph{sp} orbital for the $\sigma$
bond, one \emph{sp} orbital for its lone pair, and the two remaining \ce{2p}
orbitals (at right angles) to form the two $\pi$ bonds. So the lone pair is in an
\emph{sp} orbital.}

\why{A / C:}{give the lone pair as being in a \ce{2p} orbital. After
hybridisation the lone pair occupies the unused \emph{sp} orbital, not a pure
\ce{2p} orbital.}

\why{D:}{says the $\sigma$ bond is formed by \emph{s}--\emph{p} overlap. In
\ce{N2} both atoms are \emph{sp} hybridised, so the overlap is
\emph{sp}--\emph{sp}.}

\mcq{A2}{A}{2025 FM Paper 12 Q16}{%
Count the $\pi$ bonds from the multiple bonds:
\ce{CO2} (\ce{O=C=O}) has \textbf{two} C=O double bonds $\rightarrow$ 2 $\pi$.
\ce{HCN} (\ce{H-C#N}) has one \ce{C#N} triple bond $\rightarrow$ 2 $\pi$.
Propene (\ce{CH3-CH=CH2}) has only one C=C $\rightarrow$ 1 $\pi$.
So \ce{CO2} and \ce{HCN} are the two molecules with two $\pi$ bonds.}

\why{B / C:}{both include propene, which has only one $\pi$ bond.}

\why{D:}{omits \ce{CO2}, which also has two $\pi$ bonds.}

\mcq{A3}{C}{2022 ON Paper 11 Q5}{%
\ce{H-C#C-C(CH3)=CH(CH3)} --- count every single bond as one $\sigma$, and remember
that each double/triple bond contributes only \textbf{one} $\sigma$:
\ce{H-C} (1), \ce{C#C} (1 $\sigma$), \ce{C-C} (1), the C=C (1 $\sigma$),
\ce{C-H} on the double-bonded CH (1), and the two methyl groups
(\ce{CH3}: 1 C--C + 3 C--H each = 4 + 4). Total $= 1+1+1+1+1+4+4 = 13$.}

\why{A / B:}{these miss some C--H $\sigma$ bonds (a $\sigma$ bond exists for
every single line and for every C--H bond).}

\why{D:}{16 counts the $\pi$ bonds as if they were $\sigma$ bonds --- the two
$\pi$ bonds of the alkyne plus the one $\pi$ bond of the alkene must not be
counted.}

\mcq{A4}{C}{2024 MJ Paper 11 Q26}{%
X is \ce{NC-CH=CH2}. Count the $\sigma$ bonds: the C--C bond (1), the $\sigma$
bond inside \ce{C#N} (1), the $\sigma$ bond inside C=C (1), one C--H on the CH (1) and
the two C--H bonds on the \ce{CH2} (2). Total $= 6$ $\sigma$ bonds.
The \ce{C#N} carbon is \emph{sp}; the two C=C carbons are \emph{sp}\textsuperscript{2}. There is no
\ce{CH3} group, so there are no \emph{sp}\textsuperscript{3} carbons. Row C.}

\why{A / B:}{give a total of 4 $\sigma$ bonds --- the $\sigma$ bonds inside the
\ce{C#N} and C=C are omitted.}

\why{D:}{has the correct $\sigma$ count but wrongly includes \emph{sp}\textsuperscript{3} carbon;
X has no saturated (four-single-bond) carbon atom.}

\mcq{A5}{D}{2025 FM Paper 12 Q39}{%
but-2-ene is \ce{CH3-CH=CH-CH3}. Only the two terminal \ce{CH3} carbons are
\emph{sp}\textsuperscript{3} (the two C=C carbons are \emph{sp}\textsuperscript{2}). Each \emph{sp}\textsuperscript{3} carbon has
\textbf{4} \emph{sp}\textsuperscript{3} hybrid orbitals, all used for bonding (3 C--H and 1 C--C),
so the total is $2 \times 4 = 8$.}

\why{A / B / C:}{these count \emph{sp}\textsuperscript{3} carbon atoms or the orbitals of only one
carbon. The question asks for the total number of \emph{sp}\textsuperscript{3} orbitals.}

\mcq{A6}{B}{2025 ON Paper 11 Q6}{%
\ce{NH4CN} is made of \ce{NH4+} and \ce{CN-} ions.
\ce{NH4+}: one of its four N--H bonds is a \emph{coordinate} bond (the lone pair
on N in \ce{NH3} is donated to \ce{H+}); the other three are ordinary covalent
bonds. So \textbf{1} coordinate bond.
\ce{CN-}: the \ce{C#N} triple bond contains \textbf{2} $\pi$ bonds.
Therefore 1 coordinate bond and 2 $\pi$ bonds.}

\why{A / C:}{say there are 0 coordinate bonds --- but the \ce{NH4+} ion contains
one dative bond.}

\why{D:}{says 3 $\pi$ bonds; a \ce{C#N} triple bond has only two $\pi$ bonds (and one
$\sigma$).}

\mcq{A7}{C}{2023 ON Paper 12 Q5}{%
In \ce{NH3} the nitrogen is \emph{sp}\textsuperscript{3} hybridised (3 bond pairs + 1 lone pair).
When it donates its lone pair to \ce{H+} to form \ce{NH4+}, nitrogen is still
surrounded by four electron pairs and is still \emph{sp}\textsuperscript{3}. The hybridisation
therefore does \emph{not} change.}

\why{A:}{\ce{NH4+} is a regular tetrahedral ion, so its dipole moment is
\emph{zero}; \ce{NH3} has a significant dipole moment. The dipole moment
decreases, not increases.}

\why{B:}{the H--N--H bond angle \emph{increases} from about
107\textdegree\ (\ce{NH3}) to 109.5\textdegree\ (\ce{NH4+}), because the lone pair
is replaced by a bonding pair.}

\why{D:}{there is no electron transfer from N to Cl. A \emph{proton} (\ce{H+}) is
transferred from \ce{HCl} to \ce{NH3}; the lone pair on N is \emph{shared}, not
transferred.}

\mcq{A8}{B}{2022 MJ Paper 13 Q9}{%
The average C--H bond energy in methane is defined as the enthalpy change when
\textbf{1 mol of C--H bonds} is broken, averaged over the four (non-equivalent)
bonds. Since methane has four C--H bonds, the equation must break 1 mol of C--H
bonds in total, i.e.\ $\tfrac14$ mol of \ce{CH4} to give $\tfrac14$ mol of
\ce{C(g)} and 1 mol of \ce{H(g)} --- all species \emph{gaseous} and the C--H bonds
broken (atoms formed).}

\why{A:}{this is the reverse process (bond \emph{formation}) and would give the
negative of the bond energy.}

\why{C:}{this is the \emph{total} atomisation of one mole of methane (breaking
\emph{all four} C--H bonds), which is $4\times$ the average C--H bond energy.}

\why{D:}{this breaks only the \emph{first} C--H bond (to give \ce{CH3} and H); it
is not an average over all four bonds.}

\newpage

\section*{Section B --- Mark scheme clips}

\msblock{B1(a)(ii)  [4]}{cie-9701-2024-mj-p23-q1a-ii-ms.png}

\msblock{B2(b)(i)  [1]}{cie-9701-2024-mj-p21-q5b-i-ms.png}

\msblock{B3(b)(ii)  [2]}{cie-9701-2024-mj-p21-q5b-ii-ms.png}

\msblock{B4(c)(iii)  [2]}{cie-9701-2022-on-p22-q1c-iii-ms.png}

\msblock{B5(b)  [1]}{cie-9701-2021-mj-p21-q2b-ms.png}

\msblock{B6(d)(ii)  [2]}{cie-9701-2025-mj-p24-q3d-ii-ms.png}

\end{document}
